% !TeX program = lualatex
% Compile twice: lualatex 117.tex
% All references and prose are embedded; no BibTeX database or external inputs.
\documentclass[11pt,a4paper]{article}
\usepackage[margin=24mm,headheight=14pt]{geometry}
\usepackage{fontspec}
\setmainfont{Latin Modern Roman}
\setsansfont{Latin Modern Sans}
\setmonofont{Latin Modern Mono}
\usepackage{amsmath,amssymb,mathtools}
\usepackage{unicode-math}
\setmathfont{Latin Modern Math}
\usepackage{microtype}
\usepackage{enumitem,xurl,needspace}
\usepackage[dvipsnames]{xcolor}
\usepackage{hyperref}
\hypersetup{unicode=true,colorlinks=true,linkcolor=MidnightBlue,urlcolor=MidnightBlue,
 pdftitle={Research attempt: Problem 117},pdfauthor={Research notebook},bookmarksnumbered=true}
\usepackage{bookmark}
\usepackage{tocloft}
\setlength{\cftsecnumwidth}{3em}
\cftsetpnumwidth{2.5em}
\cftsetrmarg{4em}
\usepackage{titlesec}
\titleformat{\section}{\Large\bfseries\raggedright}{\thesection}{0.7em}{}
\usepackage{fancyhdr}
\pagestyle{fancy}\fancyhf{}
\fancyhead[L]{\small\sffamily Open Problems Math | Research notebook}
\fancyhead[R]{\small Problem 117 | 27 September 2026}
\fancyfoot[C]{\thepage}
\setlength{\parindent}{0pt}
\setlength{\parskip}{3pt plus 1pt}
\setlength{\emergencystretch}{3em}
\setcounter{tocdepth}{1}
\setcounter{secnumdepth}{1}
\setlist[itemize]{leftmargin=3em,itemsep=5pt,topsep=4pt,parsep=0pt}
\allowdisplaybreaks
\newcommand{\N}{\mathbb{N}}
\newcommand{\Z}{\mathbb{Z}}
\newcommand{\Q}{\mathbb{Q}}
\newcommand{\R}{\mathbb{R}}
\newcommand{\C}{\mathbb{C}}
\newcommand{\F}{\mathbb{F}}
\newcommand{\A}{\mathbb{A}}
\newcommand{\PP}{\mathbb P}
\newcommand{\E}{\mathbb{E}}
\newcommand{\eps}{\varepsilon}
\newcommand{\dd}{\,\mathrm d}
\newcommand{\sref}[2]{\hyperlink{source:#1:#2}{[S#2]}}
\newcommand{\eref}[2]{\hyperlink{extra:#1:#2}{[E#2]}}
% Turn the next switch off for a shorter reading copy. The quotations preserve
% every exactTarget field, including qualifications omitted from a short summary.
\newif\ifshowcatalogue
\showcataloguetrue
\newcommand{\cataloguescope}[1]{\ifshowcatalogue
 \paragraph{Exact scope in the supplied catalogue.}
 \begin{quote}\small #1\end{quote}\fi}
\usepackage{etoolbox}
\AtBeginEnvironment{itemize}{\small}
\numberwithin{equation}{section}
% Standalone export. Original section 117 preserved verbatim outside marked additions.
\begin{document}
\setcounter{section}{116}
% ================================================================
% problemId: problem.prove-the-bkl-locality-conjecture-in-the-general-inhomogeneous-context
% source releaseRank: 117; source releaseStatus: open
% exactTarget SHA-256: 34fe91bf73eb6698c705de8cbdaf273ac7f0e9536b11919fd14489986cc792dc
\clearpage
\section{\texorpdfstring{Prove the BKL locality conjecture in the general inhomogeneous context}{Prove the BKL locality conjecture in the general inhomogeneous context}}\label{117}
\subsection{Definitions and mathematical statement}
Near a spacelike singularity, a Kasner epoch has a schematic local metric
\[
 ds^2\simeq-dt^2+\sum_{i=1}^3t^{2p_i}(\omega^i)^2,\qquad
 \sum_i p_i=\sum_i p_i^2=1.
\]
Vacuum Bianchi IX (Mixmaster) behavior consists of successive anisotropic epochs and transitions in a spatially homogeneous model. BKL locality predicts that, for generic inhomogeneous collapse, time evolution asymptotically dominates coupling between distinct spatial points, with corresponding oscillatory local behavior. The catalogue does not specify a data topology or measure defining ``generic,'' a gauge, or a convergence mode. Consequently this is an asymptotic geometric program, not a fully quantified theorem determined by the single sentence alone.
\par\smallskip\textit{Formulation and concept references:} \sref{117}{1}.
\cataloguescope{Prove that the singularity in generic inhomogeneous gravitational collapse is spacelike, local, and oscillatory, and that the evolution at each spatial point approaches vacuum homogeneous Bianchi IX Mixmaster behavior.}
\subsection{Short English statement}
Near a generic gravitational singularity, does each spatial point evolve almost independently through oscillatory Mixmaster-like dynamics?
% BEGIN RESEARCH ADDITION 117
\subsection{Research attempt}

\paragraph{Current frontier and formulation audit.}
Section 3.1 of the supplied 2025 review discusses oscillatory vacuum models,
velocity-dominated special cases, numerical evidence, and spatial spikes.
It does not supply a topology defining generic initial data or a norm in
which arbitrary inhomogeneous solutions shadow Bianchi IX evolution.
The vacuum and matter settings must also be distinguished. Thus the original
editorial warning is substantive, not a missing notational convention.
No choice of a convenient meaning of ``generic'' below is substituted for
the catalogue's target. \sref{117}{1}

Ringstr\"om proves an attractor theorem for spatially homogeneous Bianchi IX
solutions; this does not estimate spatial derivatives in an inhomogeneous
Einstein evolution. Andersson--Rendall construct analytic, nonsymmetric,
quiescent singularities with scalar-field or stiff-fluid matter. Those are
important warnings about matter dependence, not counterexamples to a
properly specified generic \emph{vacuum} assertion.
\eref{117}{1}\eref{117}{2}

\paragraph{Attack route.}
Try to separate three claims often compressed into the word locality:
shrinking causal neighborhoods, small spatial-derivative terms in the field
equations, and shadowing by an autonomous homogeneous trajectory. The first
admits an explicit integrated criterion; the second requires derivative
estimates that the criterion cannot provide; the third is sensitive to
instability near exceptional epochs. The calculations below deliberately
stress-test all three implications. They are auxiliary statements about
specified models, not a formulation of the entire BKL conjecture.

\paragraph{Attempt: an integrated horizon criterion.}
Consider the diagonal, spatially homogeneous \emph{kinematic} metric
\begin{equation}\label{117:metric}
 ds^2=-dt^2+\sum_{i=1}^3 a_i(t)^2(dx^i)^2,
 \qquad 0<t\le1,
\end{equation}
and put $\tau=-\log t$ and
\[
 a_i(e^{-\tau})=\exp\!\left(-\int_0^\tau p_i(s)\,ds\right),
 \qquad
 D_i(s)=\int_0^s(1-p_i(r))\,dr.
\]
We do not yet impose the vacuum evolution equations. A causal curve,
parametrized by $t$, satisfies
$\sum_i a_i(t)^2(dx^i/dt)^2\le1$. Consequently its displacement in direction
$i$, between two times tending to zero and a time $e^{-\tau}$, is bounded by
\begin{equation}\label{117:horizon}
 d_i(\tau):=\int_0^{e^{-\tau}}\frac{dt}{a_i(t)}
       =\int_\tau^\infty e^{-D_i(s)}\,ds.
\end{equation}
Thus finiteness of all three integrals gives shrinking coordinate causal
neighborhoods, with coordinate diameter at most $2\sum_i d_i(\tau)$.
The shrinking follows from convergence of the improper integrals, not from
an interchange of an uncontrolled limit with integration.

For a fixed Kasner exponent $p_i<1$, this is the familiar exact expression
\[
 d_i(\tau)=\frac{e^{-(1-p_i)\tau}}{1-p_i}.
\]
A uniform bound $1-p_i(s)\ge\delta>0$ is sufficient but unnecessarily strong.
The following weaker integrated estimate allows arbitrarily close approaches
to a Taub exponent $p_i=1$.

\textbf{Lemma (logarithmic Taub-excursion budget).}
If, for all sufficiently large $s$,
\begin{equation}\label{117:budget}
 D_i(s)\ge (1+\delta)\log(s+2)-C,
 \qquad \delta>0,
\end{equation}
then for all sufficiently large $\tau$,
\begin{equation}\label{117:tail}
 d_i(\tau)\le\frac{e^C}{\delta}(\tau+2)^{-\delta}.
\end{equation}
Indeed, insert \eqref{117:budget} into \eqref{117:horizon} and integrate
$(s+2)^{-1-\delta}$. This proves the lemma, including its constants.
A bound only of the form $D_i(s)\ge\log(s+2)-C$ does not guarantee convergence:
the comparison integral can be harmonic. $\square$

This calculation isolates a potentially useful quantity: accumulated time
away from a Taub direction, not a pointwise lower bound on every epoch.
For actual rotating, inhomogeneous frames, one would need corresponding
bounds on the full frame coefficients along causal curves. Equation
\eqref{117:horizon} alone does not justify that extension.

\paragraph{An exact test that defeats a pointwise shortcut.}
Let $\epsilon(\tau)=1/(\tau+4)$ and define
\begin{equation}\label{117:exponents}
 p_1=1-\epsilon,\qquad
 p_2=\frac{\epsilon+\sqrt{4\epsilon-3\epsilon^2}}2,\qquad
 p_3=\frac{\epsilon-\sqrt{4\epsilon-3\epsilon^2}}2.
\end{equation}
At every time,
\[
 \sum_i p_i=1,\qquad \sum_i p_i^2=1,\qquad p_i<1.
\]
Nevertheless,
\[
 D_1(\tau)=\log\frac{\tau+4}{4},\qquad
 d_1(\tau)=\int_\tau^\infty\frac{4\,ds}{s+4}=\infty.
\]
Thus instantaneous membership in the non-Taub part of the Kasner circle
does not imply a shrinking horizon. The amount of time spent increasingly
near its Taub point matters.

Crucially, this is \emph{not} a vacuum counterexample. The Einstein equations
can be checked rather than presumed. For \eqref{117:metric}, write
$H_i=\dot a_i/a_i=p_i(\tau)/t$. Then
\[
 \dot H_i=\frac{-p_i'(\tau)-p_i(\tau)}{t^2},
 \quad
 R_{00}=-\sum_i(\dot H_i+H_i^2)=0,
 \quad
 \frac{R_{ii}}{a_i^2}=\dot H_i+H_i\sum_jH_j
                  =-\frac{p_i'(\tau)}{t^2}.
\]
Since the $p_i$ in \eqref{117:exponents} are not constant, the spatial vacuum
equations fail. Even though the dimensionless residuals
$t^2R_{ii}/a_i^2=-p_i'$ tend to zero, approximate satisfaction is not exact
satisfaction. The example invalidates a proposed \emph{kinematic inference},
not the conjecture about genuine Einstein solutions.

\paragraph{Why a shrinking coefficient is not a PDE estimate.}
In a schematic normalized evolution
\begin{equation}\label{117:schematic}
 \partial_\tau U=F(U)+\sum_a E_a^{\ i}(\tau,x)\partial_iU+\mathcal R,
\end{equation}
one cannot discard the derivative term merely because $E_a^{\ i}\to0$.
For instance, for any $\delta>0$,
\[
 E(\tau)=e^{-\delta\tau},\qquad
 U(\tau,x)=\sin(e^{\delta\tau}x)
 \quad\Longrightarrow\quad
 E\partial_xU=\cos(e^{\delta\tau}x).
\]
The coefficient vanishes while the weighted derivative has supremum norm one.
This is another elementary test function, not an Einstein solution. It
explains the exact additional estimate a proof needs, for example
$\|E_a^{\ i}\partial_iU\|\to0$ in a specified norm on a specified set of spatial
points. Spikes are especially relevant when assessing which such norm and
exceptional set could be appropriate. \sref{117}{1}

\paragraph{A shadowing stress test.}
Even small or integrable remainders in \eqref{117:schematic} need not give
tracking by the homogeneous solution with the same initial state. On an
interval on which $F$ has Lipschitz bound $L(\tau)$, subtraction from
$z'=F(z)$ and Gronwall give only
\begin{equation}\label{117:gronwall}
 |U(\tau)-z(\tau)|\le
 e^{\int_{\tau_0}^\tau L(r)dr}|U(\tau_0)-z(\tau_0)|
 +\int_{\tau_0}^\tau e^{\int_s^\tau L(r)dr}|R(s)|\,ds,
\end{equation}
where $R$ denotes the full discarded term. The exponential amplification
cannot be replaced by one without a stability argument.

There is an exact discrete illustration. For the doubling map
$T(x)=2x\pmod1$, the unperturbed orbit from zero is zero. Perturb only the
first step, setting $y_1=1/(3\cdot2^K)$, and then set $y_{j+1}=T(y_j)$.
The sum of all forcing errors is the arbitrarily small number
$1/(3\cdot2^K)$, but $y_{K+1}=1/3$, after which the orbit alternates between
$1/3$ and $2/3$. Tracking the zero orbit fails permanently. This does not
exclude shadowing by a \emph{different} orbit with adjusted initial data;
that is precisely why a correctly formulated shadowing theorem is needed.
Nor is this expanding map being identified with the complete Bianchi IX
flow.

\paragraph{Stress test against the exact target.}
None of these arguments proves that a generic singularity is spacelike.
The first lemma starts with the particular metric \eqref{117:metric}; it
cannot establish that a general Einstein solution admits this form. The
Kasner-circle example fails the spatial Einstein equations, the oscillatory
function fails to be a constructed gravitational solution, and the doubling
map is a stability test rather than a relativistic counterexample.
Homogeneous attractor theorems supply neither the missing spatial estimates
nor an automatic identification of individual inhomogeneous trajectories.
\eref{117}{1}

A viable continuation would have to specify the vacuum data class and gauge,
prove an integrated excursion bound or another causal-locality estimate for
actual solutions, control the spatial derivative terms including their
exceptional sets, and establish a suitable nonuniform shadowing statement.
Quiescent matter solutions must remain outside any claim specifically about
vacuum oscillations; suppressing that distinction would change the problem.
\eref{117}{2}

\paragraph{Outcome.}
\textbf{Outcome: Exploratory approach with unresolved gap.}

The attempt gives an explicit integrated horizon criterion and its sharp
logarithmic comparison threshold, together with three falsified shortcut
lemmas. The varying-Kasner test has its vacuum residual calculated exactly,
so it is not misreported as a counterexample. The companion script checks
the algebraic identities and the discrete instability example.

No generic inhomogeneous Einstein estimate is established. In addition to
those analytic gaps, the supplied formulation still lacks quantifiers for
genericity and convergence. The notebook preserves that limitation rather
than declaring a resolution of a newly chosen special formulation.
% END RESEARCH ADDITION 117
\subsection{Sources}
\begin{itemize}\raggedright
\item[\textnormal{[S1]}]\hypertarget{source:117:1}{} Josu C. Aurrekoetxea, Katy Clough, and Eugene A. Lim, Cosmology using numerical relativity, Section 3.1 (2025).
\url{https://link.springer.com/article/10.1007/s41114-025-00058-z}.
{\small\textit{Catalogue formulation source.}}
% BEGIN NEW REFERENCES 117
\item[\textnormal{[E1]}]\hypertarget{extra:117:1}{} Hans Ringstr\"om, \emph{The Bianchi IX attractor}, Annales Henri Poincar\'e \textbf{2} (2001), 405--500; arXiv:gr-qc/0006035. Homogeneous attractor theorem and scope of the matter cases.
\url{https://arxiv.org/abs/gr-qc/0006035}.
\item[\textnormal{[E2]}]\hypertarget{extra:117:2}{} Lars Andersson and Alan D. Rendall, \emph{Quiescent cosmological singularities}, Communications in Mathematical Physics \textbf{218} (2001), 479--511; arXiv:gr-qc/0001047. Analytic scalar-field and stiff-fluid constructions.
\url{https://arxiv.org/abs/gr-qc/0001047}.
% END NEW REFERENCES 117
\end{itemize}

\end{document}
